On dose une solution d'acide chlorhydrique de concentration $C_{A}$ par une
solution d'hydroxyde de sodium de concentration
$C_{B}=\qty{0.010}{\mol\per\L}$.

On place $\qty{10}{\mL}$ de la solution d'acide chlorhydrique dans un bécher, on
ajoute de l'eau distillée et on y fait couler, avec une burette graduée, la
solution d'hydroxyde de sodium.
\begin{enumerate}
    \item Écrire l'équation du dosage.
    \item Dresser le tableau permettant de suivre l'évolution de la
          transformation en fonction de l'avancement $x$.
    \item Définir l'équivalence. En déduire une relation, à l'équivalence,
          permettant de calculer la concentration $C_{A}$ de la solution d'acide
          chlorhydrique.
    \item On a déterminé par conductimétrie le volume de solution d'hydroxyde de
          sodium à l'équivalence~: $V_{B\mathrm{E}}=\qty{13.6}{\mL}$.

          Calculer la concentration de la solution d'acide chlorhydrique.
\end{enumerate}

